We define and analyze the IJ covariance estimate of the frequentist variance of
Bayesian posterior expectations.  The IJ covariance is computationally
appealing, since it can be computed easily from only a single set of MCMC draws,
in contrast to the bootstrap, which, in general, requires many distinct MCMC
chains.  In finite dimensional problems obeying a central limit theorem, we show
that the IJ covariance is consistent, since it consistently estimates the
sandwich covariance.  Our work rigorously connects Bayesian posterior inference
with a large and growing literature on applications of the empirical influence
function.

We caution that our proofs are asymptotic, based on finite-dimensional
posteriors, and unfortunately come with no accompanying guarantees for any
particular dataset.  Based on our experience, we are confident recommending the
IJ covariance as a rough but computationally efficient check for worryingly high
levels of frequentist sampling variance. However, in applications where precise
frequentist coverage is crucial, we recommend supplementing the IJ covariance
with more computationally expensive procedures such as the bootstrap or
conformal intervals \citep{huggins:2022:bayesbag,fong:2021:conformal}.
Furthermore, as we discuss in \secref{data_weights} above, a simple argument
suggests that the IJ covariance will be inaccurate in high-dimensional
posteriors, and the IJ should be used with extreme caution in such cases.
Forming diagnostics to alert the user when IJ covariance may be inaccurate is an
exciting and important avenue for future work.